% Part: applied-modal-logics
% Chapter: epistemic-logic
% Section: introduction

\documentclass[../../../include/open-logic-section]{subfiles}

\begin{document}

\olfileid{aml}{el}{int}

\olsection{Pambuka}

Kaya logika modal ngrembug \emph{proposisi modal} lan sesambungan
akibat logis ing antarane, logika epistemik ngrembug
\emph{proposisi epistemik} lan sesambungan akibat logis ing antarane.
Tinimbang napsirake operator modal minangka kemungkinan lan
keniscayan, konektif unèr kasebut ditafsirake kanthi cara epistemik
utawa doksastik kanggo model kawruh lan kapercayan. Contone, kita
bisa kepengin medharake pratelan kaya ing ngisor iki:

\begin{enumerate}
\item Richard ngerti yen Calgary ana ing Alberta.
\item Audrey nganggep mungkin ana asu ing sofa.
\item Richard ngerti yen Audrey ngerti yen kelase dianakake saben dina Selasa.
\item Saben wong ngerti yen setaun ana 12 sasi.
\end{enumerate}
Asal-usul logika epistemik kontemporer kerep dilacak menyang buku
\emph{Knowledge and Belief} anggitane Jaakko Hintikka taun 1962.
Buku iku ditulis nalika semantik jagad mungkin saya akeh digunakake
ing logika. Sejatine, logika epistemik nggunakake akeh piranti
semantik kang padha karo logika modal liyane, nanging ditafsirake
kanthi cara kang beda. Owah-owahan utamane ana ing makna
\emph{relasi aksesibilitas}. Ing logika epistemik, relasi iku makili
sawijining wujud \emph{kemungkinan epistemik}. Kita bakal weruh
yen gagasan epistemik kang dimodelake mengaruhi sarat kang arep
ditrapake marang relasi aksesibilitas. Kita uga bakal nliti apa kang
kedadeyan marang teori korespondensi yen diwenehi tafsir epistemik.
Conto ing ndhuwur nyebut rong agen, Richard lan Audrey, sarta
sesambungan antarane bab kang dingerteni saben wong. Logika
epistemik kang bakal kita rembug iku logika multiagen, kang bisa
medharake bab kaya mangkono. Kosok baline, logika epistemik agen
tunggal mung ngrembug bab kang dingerteni utawa dipercaya dening
siji individu.

\end{document}
